\chapter{Bar Euler products and Borcherds denominators}
\label{ch:bar-denominators}

\section{The denominator product}

Let
\[
  \mathfrak g = \mathfrak h \oplus
  \bigoplus_{\alpha\in\Gamma\setminus\{0\}} \mathfrak g_\alpha
\]
be a root-graded Lie superalgebra, $\Gamma_+\subset\Gamma$ a positive
monoid, and $\mathfrak n_+=\bigoplus_{\alpha\in\Gamma_+}\mathfrak
g_\alpha$. Assume every root space is finite dimensional and every degree
admits finitely many decompositions in the chosen $\Gamma_+$-adic
completion. Internal PBW and Theorem~\ref{thm:bar-chevalley} give
\[
  \Barperp\bigl(U(\mathfrak n_+)\bigr)
  \;\simeq\; \CE_*(\mathfrak n_+)
  \;=\; \Sym\bigl(\mathfrak n_+[1]\bigr).
\]

\begin{theorem}[Bar denominator product]
\label{thm:bar-denominator}
\StatusProved{} The completed graded Euler supercharacter of the ordered
bar is
\[
  \chi_\Gamma\Bigl(\Barperp\bigl(U(\mathfrak n_+)\bigr)\Bigr)
  \;=\;
  \prod_{\alpha\in\Gamma_+}\bigl(1-e^{-\alpha}\bigr)^{\sdim\mathfrak g_\alpha}.
\]
If $\mathfrak g$ is a Borcherds--Kac--Moody superalgebra with Weyl vector
$\rho$, then $e^{-\rho}\chi_\Gamma\bigl(\Barperp U(\mathfrak n_+)\bigr)$
is the product side of its Weyl--Kac--Borcherds denominator identity.
\end{theorem}

\begin{proof}
Choose a homogeneous basis of each root space. Suspension reverses parity.
An even root vector of root $\alpha$ becomes odd in $\Sym(\mathfrak
n_+[1])$ and contributes an exterior factor $1-e^{-\alpha}$ to the Euler
supercharacter; an odd root vector becomes even and contributes a
polynomial factor $(1-e^{-\alpha})^{-1}$. Multiplying over the basis gives
the exponent $\dim\mathfrak g_{\alpha,\bar0}-\dim\mathfrak
g_{\alpha,\bar1}=\sdim\mathfrak g_\alpha$. The finite-decomposition
hypothesis makes every coefficient of the completed product a finite sum.
The prefactor $e^{-\rho}$ is the leading Weyl monomial of the denominator
identity \cite{Bor88,Kac90}.
\end{proof}

This is an object-to-character bridge with the functor named. The bracket
controls the Chevalley differential; the Euler product remembers only the
signed root spaces. It is not an identification of the bar object with an
automorphic form, and it is not a classification.

\section{The Igusa specialisation}

Let $\phi_{0,1}(\tau,z)=\sum_{n\ge0,\ \ell\in\mathbb
Z}f(n,\ell)q^ny^\ell$ be the weak Jacobi form of weight $0$ and index $1$
normalized by $f(0,0)=10$ and $f(0,\pm1)=1$. In the standard cusp chamber
its monic Borcherds product is
\[
  D_5(Z) = q^{1/2}y^{1/2}p^{1/2}
  \prod_{(n,\ell,m)>0}\bigl(1-q^ny^\ell p^m\bigr)^{f(nm,\ell)},
  \qquad \Delta_5(Z) = 64\,D_5(Z),
\]
the second equality being the theta-constant normalization.
Gritsenko--Nikulin \cite{GN97} construct a Borcherds algebra
$\mathfrak g_{\Delta_5}$ with signed root multiplicities
$\sdim(\mathfrak g_{\Delta_5})_{\alpha(n,\ell,m)}=f(nm,\ell)$ and
denominator $\den(\mathfrak g_{\Delta_5})=D_5(2Z)$.

\begin{corollary}[Igusa bar identity]
\label{cor:igusa}
\StatusProved{} For the positive nilpotent subalgebra $\mathfrak
n_{\Delta_5,+}$,
\[
  e^{-\rho}\,
  \chi_\Gamma\Bigl(\Barperp\bigl(U(\mathfrak n_{\Delta_5,+})\bigr)\Bigr)
  \;=\; D_5(2Z) \;=\; 64^{-1}\Delta_5(2Z).
\]
\end{corollary}

\begin{proof}
Insert the Gritsenko--Nikulin root supermultiplicities into
Theorem~\ref{thm:bar-denominator}. The Weyl monomial and the doubled
period variable are those of the type-II denominator chamber, and the
theta-product leading coefficient converts the monic product to
$64^{-1}\Delta_5(2Z)$.
\end{proof}

The root multiplicities are quoted input from \cite{GN97}, not a
consequence of Theorem~\ref{thm:bar-denominator}. What the theorem
supplies is the object whose Euler character the product is.

\begin{remark}[A compact realization is a separate theorem]
\label{rem:k3e}
A geometric realization of $\mathfrak g_{\Delta_5}$ on a compact
Calabi--Yau background, in particular any $\mathrm{K3}\times E$
realization, is a separate geometric statement. It does not follow from
Corollary~\ref{cor:igusa}: matching a scalar product against a
denominator identity constructs no morphism of the structured objects that
precede the character. See No-go~\ref{nogo:scalar-sequence} and
No-go~\ref{nogo:anomaly-types}.
\end{remark}
